\clearpage\pagestyle{otMain}\setlength{\otBlocEspacement}{2.5mm}\fontsize{11}{15}\selectfont\corrigesChapitre{10}{Calculer avec les décimaux}
\begin{blocLivre}{cor-10-1}\begin{corrige}{1}{Somme et différence}\demarche{Aligner les virgules, compléter les rangs.}
$12{,}70+3{,}85=\res{16{,}55}$ ; $20{,}00-4{,}35=\res{15{,}65}$ ; contrôle : $15{,}65+4{,}35=20$.
\end{corrige}\end{blocLivre}
\begin{blocLivre}{cor-10-2}\begin{corrige}{2}{Un produit posé}\demarche{Compter les chiffres après la virgule.}
$638\times7=4\,466$, donc $6{,}38\times7=\res{44{,}66}$. Estimation : $6{,}4\times7\approx45$.
\end{corrige}\end{blocLivre}
\begin{blocLivre}{cor-10-3}\begin{corrige}{3}{Par 10, 100, 1 000}\demarche{Déplacer les chiffres de rangs, pas « la virgule magique ».}
$\res{345{,}6}$ ; $\res{70}$ ; $\res{0{,}472}$ ; $\res{0{,}005}$.
\end{corrige}\end{blocLivre}
\begin{blocLivre}{cor-10-4}\begin{corrige}{4}{Doubles et moitiés}\demarche{Travailler en dixièmes ou centièmes.}
Double : $\res{6{,}9}$. Moitié de $7{,}3$ : $\res{3{,}65}$. Moitié de $0{,}9$ : $\res{0{,}45}$.
\end{corrige}\end{blocLivre}
\begin{blocLivre}{cor-10-5}\begin{corrige}{5}{Multiplier par 5 et par 50}\demarche{Passer par 10 ou 100 puis prendre la moitié.}
$4{,}6\times10=46$, moitié : $\res{23}$. $2{,}4\times100=240$, moitié : $\res{120}$.
\end{corrige}\end{blocLivre}
\begin{blocLivre}{cor-10-6}\begin{corrige}{6}{Division en ligne}\demarche{Passer le reste au rang suivant.}
25 unités : 8, reste 1 ; 15 dixièmes : 5. Quotient $\res{8{,}5}$ ; contrôle : $8{,}5\times3=25{,}5$.
\end{corrige}\end{blocLivre}
\begin{blocLivre}{cor-10-7}\begin{corrige}{7}{Diviser par 4 et par 8}\demarche{Prendre la moitié deux ou trois fois.}
$18\to9\to\res{4{,}5}$ : deux moitiés pour $:4$. Une moitié de plus : $18:8=\res{2{,}25}$.
\end{corrige}\end{blocLivre}
\begin{blocLivre}{cor-10-8}\begin{corrige}{8}{Le quotient qui s’allonge}\demarche{Poursuivre jusqu’au reste nul.}
7 unités : 1 ; 23 dixièmes : 4, reste 3 ; 30 centièmes : 6. Quotient $\res{1{,}46}$ ; contrôle : $1{,}46\times5=7{,}3$.
\end{corrige}\end{blocLivre}
\begin{blocLivre}{cor-10-9}\begin{corrige}{9}{Compléments et regroupements}\demarche{Chercher les paires qui font un compte rond.}
$4{,}35+0{,}65=5$, puis $5+2{,}8=\res{7{,}8}$. Complément : $\res{3{,}75}$.
\end{corrige}\end{blocLivre}
\begin{blocLivre}{cor-10-10}\begin{corrige}{10}{L’estimation qui sauve}\demarche{Encadrer par un produit d’entiers ronds.}
$19{,}8$ est proche de 20 et $20\times6=120$ : le seul résultat vraisemblable est $\res{118{,}8}$.
\end{corrige}\end{blocLivre}
\begin{blocLivre}{cor-10-11}\begin{corrige}{11}{Les courses du club}\demarche{Calculer chaque produit avant la somme.}
Ballons : $12{,}45\times6=74{,}70$ euros ; pompes : $4{,}20\times3=12{,}60$ euros ; total $\res{87{,}30}$ euros. Monnaie : $90-87{,}30=\res{2{,}70}$ euros.
\end{corrige}\end{blocLivre}
\begin{blocLivre}{cor-10-12}\begin{corrige}{12}{Le carrelage du couloir}\demarche{Convertir pour diviser des entiers.}
En décimètres : $72:4=\res{18}$ dalles. Changer d’unité remplace la division décimale par une division entière.
\end{corrige}\end{blocLivre}
\begin{blocLivre}{cor-10-13}\begin{corrige}{13}{Le partage du week-end}\demarche{Partager avant de soustraire.}
Part : $25{,}5:3=\res{8{,}5}$ euros. Après dépense : $8{,}5-5{,}90=\res{2{,}60}$ euros chacun.
\end{corrige}\end{blocLivre}
\begin{blocLivre}{cor-10-14}\begin{corrige}{14}{La moyenne du relais}\demarche{Comparer chaque part à la part égale.}
Part égale : $10{,}8:4=\res{2{,}7}$ km. Le coureur fait $2{,}7-2{,}5=0{,}2$ km de moins : un autre devra courir $\res{0{,}2}$ km de plus, soit $2{,}9$ km.
\end{corrige}\end{blocLivre}
\begin{blocLivre}{cor-10-15}\begin{corrige}{15}{Multiplier rapetisse ?}\demarche{Tester sur un décimal plus petit que 1.}
$0{,}3\times10=3>0{,}3$ : la première est $\res{\text{vraie}}$ pour les nombres positifs non nuls. $0{,}3:10=0{,}03<0{,}3$ : diviser par 10 rapetisse toujours un nombre positif : la seconde est $\res{\text{fausse}}$.
\end{corrige}\end{blocLivre}
